% Draft subsection for G10, to follow Proposition prop:fgprepring and to replace the paragraph
% "When $B^M\ne1$ the preparation of all pairs ..." that follows it. Director Z, 15 November 2026. Not issued.

\subsection{Four arc gains on rings: the winding of the walks at one vertex}

On the ring the gain of a walk depends on two of its statistics, and Proposition~\ref{prop:fgprepring} avoided the second by assuming $B^{M}=1$. This subsection drops that assumption in two cases: one class of winding, given a property of the one-gain construction that follows from the statements of \cite{dgone}, and two classes on rings with $M\ge8$, by a preparation of three groups. It also shows that for three or more classes the preparation of all pairs does not exist. Throughout, $\lambda(w)$ is the number of left steps of a walk $w$, $\lambda_t(w)$ the number among its first $t$ steps, and
\[
\beta=B^{M}=(bc\,a^{-2})^{M},\qquad \Delta_N=\langle\delta,\beta\rangle,\qquad \mathsf{s}=\min\{j\ge1:\beta^{j}\in\langle\delta\rangle\};
\]
$\Delta_N$ is the group of the setting because $(a/d)^{M}=\beta^{-1}\delta^{-M}$. A walk is its sequence of step directions.

\begin{lemma}[gain and winding of one walk]\label{lem:fgkappa}
A walk from $(0,R)$ of length $T$ with $\lambda$ left steps has gain $\kappa\,g_\delta(w)$ if it ends on $R$ and $\kappa\,b\delta d^{-1}g_\delta(w)$ if it ends on $L$, with $\kappa=a^{T-2\lambda}(bc)^{\lambda}$. Replacing $\lambda$ by $\lambda+M$ at fixed $T$ multiplies $\kappa$ by $\beta$. On $\mathbb{Z}_N$ two walks of one length to one vertex have numbers of left steps congruent modulo $M$.
\end{lemma}
\begin{proof}
The first statement is the computation in the proof of Lemma~\ref{lem:fgrecode}, which concerns one walk and uses neither the line nor a common $\lambda$. The second is $a^{-2M}(bc)^{M}=\beta$. The site after $T$ steps is $T-2\lambda$ modulo $N$, so $2\lambda\equiv2\lambda'\pmod N$.
\end{proof}

So at a vertex $y$ of an event of one length the walks with $\lambda=l_0+jM$ have their gains in the coset $\kappa_0\beta^{j}\langle\delta\rangle$ (times $b\delta d^{-1}$ on $L$); call $j$ the \emph{level}. Levels $j$ and $j'$ give the same coset exactly when $\mathsf{s}$ divides $j-j'$.

\begin{lemma}[winding at one vertex]\label{lem:winding}
Let $S$ be a position product event of length $T$ on $\mathbb{Z}_N$, $N\ge1$, let $y$ be a vertex and $n(l)$ the number of walks of $S$ that end at $y$ with $\lambda=l$. For $l'-l\in M\mathbb{Z}$:
\begin{enumerate}[label=(\roman*)]
\item if $l'-l\ge2M$ then $n(l)\,n(l')\le n(l+M)\,n(l'-M)$;
\item $n(l)\,n(l+2M)\le n(l+M)\bigl(n(l+M)-1\bigr)$.
\end{enumerate}
Hence the values of $\lambda$ at $y$ form a progression of difference $M$ without gaps, and every value that is neither the least nor the greatest is taken by at least two walks.
\end{lemma}
\begin{proof}
Let $w,w'\in S$ end at $y$ with $\lambda(w)=l$, $\lambda(w')=l'$, and $\Phi(t)=\lambda_t(w')-\lambda_t(w)$. Then $\Phi(0)=0$, $\Phi(T)\ge2M$, and $\Phi$ changes by at most one in a step; let $\tau$ be the least $t$ with $\Phi(t)=M$, so $0<\tau<T$. After $\tau$ steps the sites of $w$ and $w'$ differ by $2M\equiv0$. Let $h$ be the first $\tau$ steps of $w'$ followed by the last $T-\tau$ steps of $w$, and $h'$ the first $\tau$ steps of $w$ followed by the last $T-\tau$ steps of $w'$. The sites of $h$ are those of $w'$ up to $\tau$ and those of $w$ after it, so $h\in S$; its last step is that of $w$, so it ends at $y$; and $\lambda(h)=l+M$. Likewise $h'\in S$ ends at $y$ with $\lambda(h')=l'-M$. The map $(w,w')\mapsto(h,h')$ is injective, since $\lambda_t(h)-\lambda_t(h')=\Phi(t)$ for $t\le\tau$ recovers $\tau$ and exchanging the parts after $\tau$ recovers $(w,w')$. This is (i). If $l'=l+2M$ then $\lambda(h)=\lambda(h')$ and $h\ne h'$, because their first $\tau$ steps contain different numbers of left steps; this is (ii). If $l<l'$ occur with $l'-l\ge2M$ then $l+M$ occurs by (i), which gives the absence of gaps by induction, and (ii) gives $n(n-1)\ge1$ at an inner value.
\end{proof}

\begin{lemma}[the two rails of one site]\label{lem:windingrails}
In the same situation let $n_R(l)$ and $n_L(l)$ be the numbers at $(x,R)$ and $(x,L)$. If $l'-l\ge M$ then $n_L(l)\,n_R(l')\le n_L(l+M)\,n_R(l'-M)$; if $l'-l\ge2M$ then $n_R(l)\,n_L(l')\le n_R(l+M)\,n_L(l'-M)$.
\end{lemma}
\begin{proof}
The same exchange. In the first case $w$ ends with a left and $w'$ with a right step, so $\Phi(T-1)=l'-l+1\ge M+1$ and $\tau<T-1$; in the second $\Phi(T-1)=l'-l-1\ge M$ and $\tau\le T-1$. In both $\tau<T$, so $h$ and $h'$ end where $w$ and $w'$ end.
\end{proof}

\begin{corollary}[no preparation of all pairs for three classes]\label{cor:nowind}
On no ring, for no position product event $S$ and no vertex $y$, does $m_S(y)$ consist of exactly three walks with gains in three different cosets of $\langle\delta\rangle$. If $\mathsf{s}\ge3$, then for $\kappa\in H$ the multiset $e_\kappa+e_{\kappa\beta}+e_{\kappa\beta^{2}}\in\mathcal{M}(\kappa\Delta_N)$ is prepared at no vertex. Whatever $\mathsf{s}$, a walk on $L$ at level $j$ together with a walk on $R$ at level $j+1$ and no other walk never occurs at one site of a position product event.
\end{corollary}
\begin{proof}
Three walks in three cosets have three levels; by Lemma~\ref{lem:winding} an inner level carries two walks. The last statement is Lemma~\ref{lem:windingrails} with $n_L(l+M)=0$.
\end{proof}

So for $\mathsf{s}\ge3$ the hypothesis under which Theorem~\ref{thm:fgsolve} with $\Delta=\Delta_N$ would classify the rules, the preparation of all pairs in $\mathcal{M}(\kappa\Delta_N)\times\mathcal{M}(\kappa bd^{-1}\Delta_N)$, is false; for $\mathsf{s}=2$ it is not obtained, and the last statement of the corollary excludes one pattern of levels. What can be prepared is this.

\begin{proposition}[one level]\label{prop:fgonelevel}
Let $N\ge3$ and suppose that $D_\delta$ on $\mathbb{Z}_N$ has the property $P^{+}_\delta$: for every site $x$ and all $U',V'\in\mathcal{M}(\langle\delta\rangle)$ there is a position product event with $m^\delta(x,R)=U'$, $m^\delta(x,L)=V'$ in which all walks that end at the site $x$ have one number of left steps. Then for arbitrary four gains, every $\kappa\in H$, $U\in\mathcal{M}(\kappa\langle\delta\rangle)$ and $V\in\mathcal{M}(\kappa bd^{-1}\langle\delta\rangle)$ are prepared at one site; a site-independent rule consistent with the factor $2$ is, on the multisets whose gains lie in one coset $\kappa\langle\delta\rangle$ or $\kappa b\langle\delta\rangle$, $\kappa\in H$, the function $\alpha|m|$ for $r$ odd and $\alpha|m|+\gamma|\sigma_\chi(m)|^{2}$, $\chi(\delta)=-1$, for $r$ even; and if $\mathsf{s}=1$ the conclusion of Theorem~\ref{thm:fgline}(i) holds.
\end{proposition}
\begin{proof}
As Proposition~\ref{prop:fgprepring}, with the common number of left steps in place of $B^{M}=1$: by Lemma~\ref{lem:fgkappa} all walks of the event at the site have one $\kappa=a^{x_0+Ni}(bc)^{\lambda_1}$; the prefixes ``left, right'' and ``$N$ steps right'' are single walks of $D_\delta$-gain $1$ back to $(0,R)$ for $N\ge3$, shift every $\lambda$ by $1$ and by $0$, and multiply $\kappa$ by $bc$ and $a^{N}$. Then Lemma~\ref{lem:fgonesite} and Theorem~\ref{thm:fgsolve} with $\Delta=\langle\delta\rangle$. If $\mathsf{s}=1$ every multiset that occurs lies in one coset, by Lemma~\ref{lem:fgkappa}.
\end{proof}

On rings of $1$, $2$ and $4$ sites $P^{+}_\delta$ is false for any gains: there both neighbours of a site are all the sites of the other parity, so every position product event leaves $0$ or a power of $2$ walks at a site. The construction of Lemma~\ref{lem:prepq} has the property in every replay made: $8{,}160$ requests on rings of $7$ to $24$ sites with $M\ge6$ and orders $1$ to $4$, each met with one number of left steps at the target. The reason is in the description of \cite{dgone}. There a walk with $\lambda$ left steps after $t$ steps is recorded by its \emph{deficit} $\lambda$, a mask is a set of deficits read modulo $M$, and each stage gives, for each wire, the integer deficits of its walks relative to a common origin: a table for the stage of five ticks, masks and step words for the short stages. In every stage the two candidates of an occupied entry differ from every retained entry of the next tick by at most five, by the statements of \cite{dgone}; so for $M\ge6$ a candidate is congruent to a retained entry modulo $M$ only if it equals it, and the admitted walks are those that the stage lists, with the integer deficits it lists. Hence all walks of one stock have one number of left steps, and at the last stage the stock on wire $2$ arrives on $L$ by one right and two left steps and the stock on wire $4$ on $R$ by three right steps, both at relative deficit $4$. So $P^{+}_\delta$ holds for $M\ge6$ as far as the statements of \cite{dgone} about its stages hold; that note checks the words of its five-tick stage in Lean and states the short stages by test, and they are not reproved here. The replays agree.

\begin{proposition}[three groups on two levels]\label{prop:threegroups}
Let $M=N/\gcd(N,2)\ge8$ and let $\delta$ have finite order $q\ge1$. For every site $x$ and all $G_1,G_2,G_3\in\mathcal{M}(\langle\delta\rangle)$ there are position product events in $D_\delta$ on $\mathbb{Z}_N$ in which the walks that end at the site $x$ are exactly
\begin{enumerate}[label=(\roman*)]
\item on $R$ the multiset $G_2$ at one level $l$; on $L$ the multiset $G_3$ at the level $l$ and $G_1$ at the level $l+M$;
\item on $R$ the multiset $G_3$ at the level $l$ and $G_1$ at $l+M$; on $L$ the multiset $G_2$ at $l+M$;
\end{enumerate}
and no walk ends at another site. Here ``at the level $l$'' means with $\lambda=l$.
\end{proposition}

\begin{proof}
\emph{Frame.} As in \cite{dgone}, record a walk by its deficit, the number of its left steps; after $t$ steps it is at the site $t-2e$ if its deficit is $e$, so a set $E$ of deficits read modulo $M$ is the mask $\{t-2e:e\in E\}$, and a walk at deficit $e$ has the two candidates $e$ (right step) and $e+1$ (left step), of which those in $E$ modulo $M$ survive. A left step multiplies the gain by $\delta$ exactly when it follows a left step. Deficits are written relative to that of the \emph{source}, a single walk of gain $\delta^{\rho}$; a \emph{stock} is a set of walks at one relative deficit $k\in\{2,4,6\}$. At the boundaries between stages every walk has arrived by a right step. In the advances, deposits and the runs made before the source is removed the relative deficits lie in $\{0,\dots,9\}$, and the checks use that for $M\ge8$ the residues of $0,\dots,7$ are distinct, that of $8$ is $0$ or $8$ and that of $9$ is $0$, $1$ or $9$. In the runs to $M$, the merge and what follows, a stock passes the entries $7,\dots,M+2$, and the checks compare the residues of entries up to $M+3$ with those of the stocks at $2$ and $4$.

\emph{Stages.} In each line the mask is the set of the entries named, and ``refused'' lists the candidates that must not lie in it.
\begin{itemize}
\item Advance (four ticks): the source goes $0,1,2,2,2$ by $LLRR$ and every stock $k,k+1,k+1,k+2,k+2$ by $LRLR$. Refused: at tick 1 the entries $0$ and $k$ against the mask of $1$ and the $k+1$; at tick 2 the entries $1$ and $k+2\in\{4,6,8\}$ against $2$ and the $k+1\in\{3,5,7\}$; at tick 3 the entries $3$ and $k+1$ against $2$ and the $k+2$; at tick 4 the entries $3$ and $k+3\in\{5,7,9\}$ against $2$ and the $k+2\in\{4,6,8\}$. The source gains one factor $\delta$, the stocks none, all deficits grow by two, and the relative deficits are unchanged.
\item Deposit (three ticks): masks $\{0,1,2\}$, $\{0,2\}$, $\{0,2\}$, each with the stocks. The source stays at $0$ by $RRR$; a copy goes $1,2,2$ by $LLR$, with one factor $\delta$, and is then a walk of the stock at $2$. Refused: the $k+1$ at every tick, $1$ at ticks 2 and 3, $3$ at tick 3.
\item Run of a stock from $k$ to $k+1$ (two ticks, $LR$): the mask is $k+1$ with the other occupied entries, twice. Refused: $k$, then $k+2$, and $j+1$ for every other occupied entry $j$. This needs that no other packet is at $k$, $k+1$ or $k+2$ modulo $M$ and that no two others are adjacent modulo $M$. Gains are unchanged.
\item Removal of the source (one tick): the mask is the set of stocks; refused are $0$, $1$ and the $k+1$.
\item Merge of a stock at $0$ into the stock at $2$ (two ticks and an idle tick): masks $\{1\}$ and $\{2\}$ with the other stocks, then the stocks. The stock goes $1,2$ by $LL$, with one factor $\delta$; refused are $0$, then $1$, and $j+1$ for the others. On the idle tick every walk takes a right step; refused are $j+1$ for every occupied entry $j$, the merged stock included ($M+3\equiv3$).
\item Coast: the mask is the set of occupied entries; every walk takes a right step; refused are $j+1$ for every occupied entry $j$.
\item End, with stocks at $p$ and $p-2$ (two ticks): masks $\{p-1,p\}$, $\{p\}$. The stock at $p$ takes $RR$ and ends on $R$; the stock at $p-2$ takes $LL$, with one factor $\delta$, and ends on $L$ of the same site; refused are $p-2$, $p+1$, then $p-1$, $p+1$.
\end{itemize}
In each case the refused candidates are not in the mask modulo $M$: these are the comparisons of residues named above. $M\ge8$ is used in the run of the first group from $2$ to $6$, whose last refused candidate is $7$ and must not be the source at $0$, in the run of the second group, where the source at $0$ and the stock at $6$ must not be adjacent, and in the advances and deposits made with a stock at $6$, whose candidates reach $7$, $8$ and $9$. For $M=7$ the first of these fails.

\emph{Schedule.} Call the groups $A=G_1$, $B=G_2$, $C=G_3$. Deposit the walks of $A$, one deposit for each, and run the stock from $2$ to $6$; deposit $B$ and run it from $2$ to $4$; deposit $C$ at $2$. Before each deposit the source is advanced until its exponent is the one wanted, the walks of a group being taken in the cyclic order of their exponents, so at most $q-1$ advances are made in a group. Remove the source. Run $A$ from $6$ to $M\equiv0$; the entries $7,\dots,M-1,0,1$ are free of the stocks at $2$ and $4$. Merge $A$ into the stock at $2$. For mode (i) end with $p=4$: on $R$ the group $B$, on $L$ the groups $C$ and $A$. For mode (ii) run $B$ from $4$ to $M\equiv0$ and end with $p=2$: on $R$ the groups $C$ and $A$, on $L$ the group $B$. Before the end, coast for as many ticks as bring the final site to $x$. An empty group is omitted with its stages; the checks only used the absence of other packets.

\emph{Levels and gains.} The number of left steps of a walk is its deficit as an integer. With $\sigma$ the deficit of the removed source, the stock $C$ is at $\sigma+2$, $B$ at $\sigma+4$ or, after its run in mode (ii), at $\sigma+M$, and $A$ at $\sigma+M+2$ after its run and merge. The end adds $2$ to the stock that ends on $L$. So in mode (i) $B$ and $C$ end with $\lambda=\sigma+4$ and $A$ with $\sigma+4+M$; in mode (ii) $C$ ends with $\sigma+2$, and $A$ and $B$ with $\sigma+2+M$. A walk deposited when the source had exponent $\rho$ carries $\delta^{\rho+1}$ and gains one further factor at a merge and one if it ends on $L$: in mode (i) the exponents are $\rho+1$ for $B$, $\rho+2$ for $C$, $\rho+3$ for $A$; in mode (ii) $\rho+1$ for $C$, $\rho+2$ for $A$ and for $B$. The advances are chosen accordingly. No walk survives outside the two stocks named, and they end at one site.
\end{proof}

For $M\le7$ this construction fails; no other was sought. The construction and its stage checks are a builder's; the proof above is the author's reading of them, tick by tick. Replays of the whole construction by the builder ($48{,}684$ requests) and by a challenger from the description alone ($14{,}892$ requests), on rings with $M$ from $8$ to $20$, met every request \cite{comp4s}.

\begin{theorem}[two classes of winding]\label{thm:fgtwo}
Let $D$ be the ring $\mathbb{Z}_N$ with $M\ge8$ and four gains in a finite abelian group with $\mathsf{s}=2$; write $\beta^{2}=\delta^{k}$ (for $r$ even the parity of $k$ is determined). Let $F_R$, $F_L$ not depend on the site and vanish at $\emptyset$, and let \eqref{eq:ppcons} hold with the factor $2$. Then on the multisets that occur $F_L=F_R\circ T_c$, $F_R$ is invariant under translations by $H$, and on $\mathcal{M}(\Delta_N)$
\[
F_R(m)=\alpha|m|+q\bigl(\varepsilon_0^{2}+\varepsilon_1^{2}\bigr)+c\,\varepsilon_0\varepsilon_1,\qquad \varepsilon_i(m)=\sum_{j}(-1)^{j}m(\beta^{i}\delta^{j}),
\]
with $q=c=0$ if $r$ is odd and $c=0$ if $k$ is odd. These are the real combinations of the count and of the squares $|\sigma_\chi|^{2}$, $\chi$ a character of $\Delta_N$ with $\chi(\delta)=-1$, one for each pair $\{\chi,\chi^{-1}\}$: for $r$ even two squares if $k$ is even and one if $k$ is odd, for $r$ odd none. Conversely each of them is consistent on every event (Theorem~\ref{thm:fgline}(ii)).
\end{theorem}
\begin{proof}
\emph{The pairs available.} By Proposition~\ref{prop:threegroups}, Lemma~\ref{lem:fgkappa} and the two prefixes of Proposition~\ref{prop:fgonelevel}, which shift all $\lambda$ equally and leave the residues $t-2\lambda$ of a later mask unchanged, there are for every $\kappa\in H$ events with $m(x,R)=U$, $m(x,L)=V$ for all $U\in\mathcal{M}(\kappa\Delta_N)$, $V\in\mathcal{M}(\kappa bd^{-1}\Delta_N)$ such that $U$ or $V$ lies in one coset of $\langle\delta\rangle$: mode (ii) gives $U=T_{\kappa'}G_3+T_{\kappa'\beta}G_1$, $V=T_{\kappa'\beta b\delta d^{-1}}G_2$, mode (i) gives $U=T_{\kappa'}G_2$, $V=T_{\kappa'b\delta d^{-1}}(G_3+T_\beta G_1)$, and $\kappa'$ may be taken in $\kappa\langle\delta\rangle$ or in $\kappa\beta\langle\delta\rangle$, since $\beta\in H$ and $\beta^{2}\in\langle\delta\rangle$. Call these pairs \emph{admissible}. By Lemma~\ref{lem:fgonesite}, \eqref{eq:fg} holds at every admissible pair.

\emph{Reduction.} Parts (i) and (ii) of Theorem~\ref{thm:fgsolve} use \eqref{eq:fg} with $V=\emptyset$ or $U=\emptyset$, which are admissible, and a change of variables that preserves admissibility. So $F_L=F_R\circ T_c$, $F_R\circ T_h=F_R$, and the restriction $f$ of $F_R$ to $\mathcal{M}(G)$, $G=\Delta_N$, is invariant under $G$ and satisfies \eqref{eq:fgQ} whenever $u$ or $v$ lies in one coset of $\langle\delta\rangle$.

\emph{Steps 1 to 5 of Theorem~\ref{thm:fgsolve}(iii).} Step~1 uses the pair $(u,E_g)$. With $l$ as there, $G_0=f-l$ is invariant, satisfies \eqref{eq:fgQ} at admissible pairs and $G_0(u+E_g)=G_0(u)$. Step~2: $G_0(W)=G_0(W+n\mathbf{1})$ extends $G_0$ to $\mathbb{Z}^{G}$; the equation persists for $u,v\in\mathbb{Z}^{G}$ with $v$ in one coset $C$ (apply it to $u+n\mathbf{1}$ and $v+nE_C$, where $E_C$ is the sum of the elements of $C$) and with $u$ in one coset (apply it to $u+nE_C$ and $v+n\mathbf{1}$). Step~3: at $(W-e_g,e_g)$ and at $(e_g,W-e_g)$ the right sides agree, so $G_0(W-e_g+e_{g\delta})=G_0(e_g+T_\delta W-e_{g\delta})=G_0(W+e_{g\delta^{-1}}-e_g)$ by invariance; replacing $W$ by $W+e_g-e_{g\delta^{-1}}$, $G_0$ is invariant under translation by $e_{h\delta^{2}}-e_h$ for every $h$. Step~4, $r$ odd: modulo the lattice $L$ of these translations and the $E_g$, $T_\delta$ is the identity and $re_g\equiv0$; the equation for $v$ in one coset reads $G_0(u+v)=G_0(u)+G_0(v)$, and since every $W$ is a sum of two such $v$, $G_0$ is additive, and $rG_0(W)=G_0(rW)=0$. Step~5, $r$ even: $P_g=e_g+e_{g\delta}$ lies in one coset and $T_\delta P_g-P_g\in L$, so $G_0(U+P_g)=G_0(U)+G_0(P_g)$; and $E_g$ is a sum of $r/2$ translates of $P_g$ congruent to it modulo $L$, so $G_0(P_g)=0$. The $P_g$ generate the kernel of $\varepsilon=(\varepsilon_0,\varepsilon_1)\colon\mathbb{Z}^{G}\to\mathbb{Z}^{2}$, so $G_0=\gamma\circ\varepsilon$, and since $\varepsilon\circ T_\delta=-\varepsilon$,
\begin{equation}\label{eq:gammaaxis}
\gamma(s+t)+\gamma(s-t)=2\gamma(s)+2\gamma(t)\qquad\text{for all }s\in\mathbb{Z}^{2}\text{ and all }t\text{ on a coordinate axis.}
\end{equation}
\emph{In place of Step~6.} On each axis \eqref{eq:gammaaxis} is the full equation on $\mathbb{Z}$, so $\gamma(n,0)=n^{2}\gamma(1,0)$ and $\gamma(0,n)=n^{2}\gamma(0,1)$. Let $\varphi(s_0,s_1)=\gamma(s_0,s_1)-\gamma(s_0,0)-\gamma(0,s_1)$. Subtracting from \eqref{eq:gammaaxis} with $t=(t_0,0)$ the same equation at $s_1=0$ gives $\varphi(s_0+t_0,s_1)+\varphi(s_0-t_0,s_1)=2\varphi(s_0,s_1)$, so $\varphi$ is affine in $s_0$; it vanishes at $s_0=0$, so $\varphi(s_0,s_1)=s_0\varphi(1,s_1)$. In the same way $\varphi(1,s_1)=s_1\varphi(1,1)$. Hence $\gamma=q_0s_0^{2}+q_1s_1^{2}+cs_0s_1$.
\emph{Invariance under $\beta$.} $\varepsilon_1(T_\beta W)=\varepsilon_0(W)$ and $\varepsilon_0(T_\beta W)=(-1)^{k}\varepsilon_1(W)$, because $\beta^{-1}=\beta\delta^{-k}$. So $q_0=q_1$ and $c=(-1)^{k}c$. Finally $l$ is additive and invariant under $G$, so $l=\alpha|m|$. A character with $\chi(\delta)=-1$ has $\chi(\beta)^{2}=(-1)^{k}$ and $\sigma_\chi=\varepsilon_0+\chi(\beta)\varepsilon_1$; for $k$ even the two squares $(\varepsilon_0\pm\varepsilon_1)^{2}$ span the forms found, for $k$ odd $|\varepsilon_0\pm i\varepsilon_1|^{2}=\varepsilon_0^{2}+\varepsilon_1^{2}$.
\end{proof}

\begin{remark}\label{rem:winding}
(a) \emph{What the theorem adds to the line.} Let $r$ be even. On the line and for $\mathsf{s}=1$ there is one square. For $\mathsf{s}=2$ and $k$ even there are two, $(\varepsilon_0+\varepsilon_1)^{2}$ and $(\varepsilon_0-\varepsilon_1)^{2}$: they add walks one winding apart with the sign $\chi(\beta)=1$ or $-1$, and a consistent rule may mix them with any real coefficients. For $k$ odd there is one square and no such freedom. The consistency requirement does not choose among the mixtures. Compare Corollary~\ref{cor:masklist}, where the quantity that position masks read is $\operatorname{Re}\chi(B)^{M}=\operatorname{Re}\chi(\beta)$. For $r$ odd only the count remains.
(b) \emph{A phase per winding that is not a rule.} For any unit $u$ the function $|\sum_w u^{\lambda(w)}\psi(w)|^{2}$ of the walks at a vertex, $\psi$ a character with $\psi(\delta)=-1$, is consistent on every event: it is the square for the carrier in which every left step carries the further gain $u$. It is a rule in the sense of this section only where it is a function of the multiset of gains. Whether it is depends on the carrier and the length. For example, on the ring of $3$ sites with $A=\mathbb{Z}_2\times\mathbb{Z}_9$, $a=(0,1)$, $b=c=0$, $d=(1,-1)$ (so $\mathsf{s}=3$), with $\psi$ the sign of the first coordinate and multisets taken up to translation by $H$ and without regard to the rail, the multiset of gains determines the function for every unit $u$ on all position product events of length at most $8$, and at length $9$ one multiset occurs with two level structures on which the functions differ unless $u^{9}=1$ \cite{comp4s}. For other gains in the same group a challenger found the first such clash at lengths $7$, $8$ or $10$.
(c) \emph{Open.} For $\mathsf{s}\ge3$: whether the count and the squares $|\sigma_\chi|^{2}$, $\chi$ a character of $\Delta_N$ with $\chi(\delta)=-1$, are the only consistent site-independent rules on the multisets that occur. They are consistent (Theorem~\ref{thm:fgline}(ii)); the proof of the converse through a preparation of all pairs is closed by Corollary~\ref{cor:nowind}. On the ring of $4$ sites with $\mathbb{Z}_2\times\mathbb{Z}_3$ and $\mathsf{s}=3$, exact ranks of the consistency equations of all position product events to length $14$ leave, on the multisets seen by length $6$, exactly the count and the two squares \cite{comp4s}; this is \textbf{COMPUTED} and supports the converse on that ring only. For $\mathsf{s}=2$: rings with $M\le7$. For $\mathsf{s}=1$: a proof of $P^{+}_\delta$ for $M\ge6$ that does not rest on \cite{dgone}, and the rings on which $P^{+}_\delta$ is false or not tested ($M\le5$).
(d) \emph{Status.} Lemmas~\ref{lem:fgkappa} to~\ref{lem:windingrails}, and Corollary~\ref{cor:nowind} are \textbf{PROVED}. Proposition~\ref{prop:threegroups} and Theorem~\ref{thm:fgtwo} are \textbf{PROVED}; the proof of the proposition is a check of finitely many small comparisons, written out above and replayed. Proposition~\ref{prop:fgonelevel} is \textbf{PROVED} given $P^{+}_\delta$, and $P^{+}_\delta$ for $M\ge6$ rests on the isolation statements of \cite{dgone} and is otherwise \textbf{COMPUTED}. The step that replaces Step~6 is checked in Lean~4 with Mathlib (\texttt{WindingAxis.lean}): an equation required only along the two axes gives a quadratic form. A challenger found Theorem~\ref{thm:fgtwo} to hold by exact ranks over the rationals on $51$ carriers with separate unknowns for the two rails and multisets of up to $4$ to $7$ gains ($r$ up to $4$). The lemma on the winding was put by the author; a builder found Lemma~\ref{lem:windingrails}, disproved the preparation of four groups it had been asked for, built Proposition~\ref{prop:threegroups} on the wires of \cite{dgone} and showed that the restricted family of pairs suffices; the proof of Theorem~\ref{thm:fgtwo} as written is the author's. The reviewers are sub-agents of the author's model family; that is not independent review.
\end{remark}
